\hypertarget{literature-survey}{%
\section{Literature Survey}\label{literature-survey}}

\hypertarget{confidence--and-budget-only-cascades}{%
\subsection{Confidence- and budget-only
cascades}\label{confidence--and-budget-only-cascades}}

BranchyNet (Teerapittayanon et al., 2016) installs side-branch
classifiers on a deep network and allows a sample to leave early as soon
as softmax entropy falls under a learned threshold. This is an idea that
has proven very fruitful to the point that early features are usually
very discriminative for easy inputs and the design is still the basis of
most cascading systems that have been developed later on. The decision
to exit the network remains a function of the input only: BranchyNet has
neither a cloud queue depth nor a network round-trip time (RTT) signal
nor an SLA budget nor a health indicator (of a running device) so it
simply cannot find out that the next layer is currently overloaded.

MSDNet (Huang et al., 2018) takes the idea further by enabling both
"anytime" and "budgeted-batch" classifications - i.e. spending a fixed
computational budget differently on easy and hardest parts of inputs
within one model, So getting rid of the necessity of having separated
networks per computation budget which had been the case before. Although
the budget itself is a design-time only constant, not a measurement
capturing the current infrastructure state. It cannot see the growth
of the cloud queue or the RTT spike just one second ago.

EdgeBERT (Tambe et al., 2021) is, by far, the most similar previously
developed system to infrastructure-aware routing and it should be
recognised for it: The use of entropy for early exit prediction is
linked to dynamic power scaling (DVFS) on per-sentence basis which
adapts the power capabilities of the chip to the difficulty level of the
input in real time, supported by actual accelerator hardware. But its
state is solely restricted to the device\textquotesingle s own chip - it
has no interface connecting it to the external signals of
infrastructures such as the status of the cloud queue or the health of
the network, because the closure of that loop has never been one of its
problems.

SPINN (Laskaridis et al., 2020) splits a single deep network across the
edge device and the cloud, running the early layers on-device and the
remainder on a cloud partner, with the exit-versus-offload decision
conditioned on both softmax entropy and a measured network connectivity
signal rather than entropy alone. This is the closest any surveyed
system comes to combining an on-device signal with a network-side one,
and it deserves the same critical treatment as its neighbours. The
connectivity signal SPINN reads is a point measurement of link quality
taken at split time, not a live cloud queue depth or a per-request SLA
deadline. It therefore still cannot tell a saturated cloud endpoint from
a lightly-loaded one. Its split point is also fixed once the network is
partitioned, rather than adapted from an outcome fed back after the
request completes --- the same limitation BranchyNet and MSDNet share.

\hypertarget{learned-single-signal-routers}{%
\subsection{Learned, single-signal
routers}\label{learned-single-signal-routers}}

RouteLLM (Ong et al., 2024) teaches a router purely from human
preference data to distinguish between a stronger and a weaker LLM by a
few data points and the system shows more than a 2x cost reduction while
also generalizing well to model pairs the router has not seen in
training. That is a significant result, being the closest deployed
comparison for a learned router. But the only context RouteLLM knows is
the prompt itself. It is not only unfamiliar but has never come across
the concept of serving-system load. In fact, from a RouteLLM
router\textquotesingle s viewpoint, a flooded cloud endpoint and a spare
one would look the same and be handled in similar ways.

The NSGA-II router (Yu et al., 2025) goes much further towards
multi-objective, infrastructure-aware routing. It uses an evolutionary
algorithm to manage and balance quality, speed, and the cost of a
response through heterogeneous edge and cloud nodes. It also keeps up
with the diversity of requests and nodes. The problem is that even
though the set of optimization objectives is the correct one, they are
only used offline to generate a Pareto-optimal policy fixed ahead of
requests, unlike a router, which must react to changes in live queue or
RTT during serving to make decisions that are request- or -queue-aware.
The PROTEUS model (Bhatti et al., 2026) is the closest in approach to
the system built in this project. It\textquotesingle s a Lagrangian-constrained RL
policy that accepts a runtime accuracy target as input, routes over a
series of LLMs, while enforcing a minimum level of accuracy.

The study demonstrated how the model could achieve accuracy up to
1.3\%-up to 4.6\% close to an oracle with cost savings of up to 89.8\%
vis-a-vis the best fixed model, a very impressive
outcome of RL applied to routing. On the downside, the state space of
PROTEUS does not contain any information about queues, latency energy
etc. The algorithm is trained to trade off precision against cost
without considering such infrastructure changes upon which the routing
in this paper is based.

GreenServ (Ziller et al., 2026) is the closest of the three learned
routers on the energy dimension: it routes each query to whichever
model in a heterogeneous pool of large language models has historically
balanced accuracy and energy use best for a query with that profile,
using a multi-armed bandit that keeps learning online rather than a
policy fixed once at training time. That online adaptivity is exactly
what BranchyNet, MSDNet and EdgeBERT lack. But the context GreenServ
conditions on is entirely query side --- task type, semantic cluster,
text complexity --- plus its own model pool's observed accuracy/energy
history. Like RouteLLM, it has no channel for cloud queue depth, network
RTT, per-request SLA remaining, or the operational health of the device
issuing the request, so two structurally identical requests are routed
identically by GreenServ regardless of whether the infrastructure behind
them is idle or saturated at that moment.

\hypertarget{cluster-aware-serving-systems}{%
\subsection{Cluster-aware serving
systems}\label{cluster-aware-serving-systems}}

Patel et al. (2024), in their paper on Splitwise, highlight that it is
the most operationally-aware system reviewed in this chapter so
far: the system performs an operation-wise (compute v. memory
bottleneck) phase separation and runs each phase on separate machines
that specialize in that phase. Splitwise explains the differences in the
phase capabilities (throughput power etc.) so well because it is a real
case of heterogeneous hardware in servers. Scheduling is at the level of
clusters and phases, but routing decisions are not accompanied by a
per-request SLA deadline check.

\hypertarget{commercial-and-open-source-landscape}{%
\subsection{Commercial and open-source
landscape}\label{commercial-and-open-source-landscape}}

The nine academic systems above are not the only relevant comparison: a mature market of "LLM gateway" products already
solves a version of the same routing problem in production, and a fair
positioning has to account for what they actually do, not only what
research papers propose. Two of the most widely deployed were checked
against their own documentation rather than their marketing pages.

\textbf{LiteLLM} (an open-source proxy self-hosted in front of 100+
model providers) documents six routing strategies for its
\texttt{Router} class: simple-shuffle (weighted pick, the recommended
production default), latency-based (lowest recent response time),
usage-based (lowest current token throughput), least-busy (fewest
in-flight calls), cost-based, and a pluggable custom-strategy interface.
It also has a cooldown mechanism that temporarily stops routing to a
deployment after a run of rate-limit or failure responses, and a
latency cache that records recent response times per deployment. This
is a genuinely reactive system: it does respond to measured latency and
failures. Every one of its strategies is nonetheless a fixed rule
evaluated against a rolling window of recent responses, so none of them
is a learned policy. None reads a self-reported edge-device health or
calibration state either, since LiteLLM routes between cloud model
endpoints rather than edge devices, leaving no analogous self-report to
read. Nor does any strategy enforce a per-request deadline before
dispatch: a request goes to whichever deployment the current rule
favours, and if that deployment is slow, the outcome is observed only
after the fact.

\textbf{Portkey}'s AI Gateway documents its load-balancing as static
weighted distribution across configured providers or models --- traffic
is split according to weights the operator sets (e.g.\ for cost
optimisation or gradual migration testing), normalised to sum to 100\%,
not according to a live queue, latency, or health signal measured at
request time. Portkey's own documentation for this feature makes no
mention of reinforcement learning, closed-loop adaptation from
outcomes, or per-request SLA enforcement. Those remain configuration
choices made by the operator up front, not decisions the gateway itself
learns or adapts.

The wider landscape (OpenRouter, Envoy AI Gateway, Requesty, and
Martian's now-discontinued standalone "intelligence router") occupies
the same territory: cost- and latency-aware routing across a pool of
model providers, marketed as adaptive but, on the evidence available in
each product's own documentation, implemented as configured rules or a
pre-trained model-performance predictor rather than a system that
learns from live infrastructure telemetry and its own routing outcomes.
None of them is edge-cloud specific in the sense this project is ---
they route between cloud model providers, not between a resource-constrained edge device and a cloud fallback --- so the edge-health
self-report, the deadman's switch, and the bidirectional trust loop
described in Section 5 do not have a commercial analogue to compare
against at all. This is independently corroborated by Li et al. (2025),
a systematic survey of edge-cloud LLM collaboration published within
weeks of this project's own evaluation work: it states plainly that it
is building "the first systematic foundation" for a unified taxonomy of
edge-cloud collaboration, and its taxonomy --- organised around task
assignment, task division, and mixture-based collaboration --- contains
no category describing a system that combines live infrastructure
signals, edge self-report, a closed feedback loop, and a learned routing
policy in one deployment. That gap in a survey published this recently
is independent evidence the gap is real, not an artefact of this
project's own literature search missing something that already exists.

\hypertarget{comparison-and-positioning}{%
\subsection{Comparison and
positioning}\label{comparison-and-positioning}}

Table~\ref{tab:comparison} summarises the eleven systems above --- nine from the
academic literature and two production LLM gateways --- against
four properties that recur through this chapter: whether the routing
signal includes anything beyond the request itself (infrastructure
awareness), whether a per-request deadline is checked before dispatch
(not just an aggregate SLA target), and whether the outcome of a routing
decision changes the state the next decision is made from (a closed
loop). No system reviewed combines all three with a per-request SLA
check. The closest, Splitwise and the NSGA-II router, are each strong on
exactly one axis and silent on the other two.

{\small\setlength{\tabcolsep}{4pt}
\begin{longtable}{p{2.3cm}p{2.9cm}p{1.5cm}p{1.5cm}p{1.3cm}p{3.2cm}}
\caption{Comparison of cascade, offloading and gateway systems surveyed in this chapter against this project's system.}\label{tab:comparison}\\
\toprule
System & Primary routing signal & Infra-aware & Per-req.\ SLA & Closed loop & Key limitation \\
\midrule
\endhead
BranchyNet (Teerapittayanon et al., 2016) & Softmax entropy at each exit & No & No & No & Exit decision is a function of the input alone, blind to what is downstream \\
MSDNet (Huang et al., 2018) & Fixed computational budget & No & No & No & Budget is a design-time constant, not a live measurement \\
EdgeBERT (Tambe et al., 2021) & Entropy-driven on-chip DVFS & Own chip only & No & No & No interface to external signals (cloud queue, network) \\
SPINN (Laskaridis et al., 2020) & Entropy plus network connectivity & Point measurement & No & No & Connectivity is a split-time reading, not a live queue or deadline \\
RouteLLM (Ong et al., 2024) & Prompt content, preference-learned & No & No & No & Cannot distinguish a saturated cloud endpoint from a spare one \\
GreenServ (Ziller et al., 2026) & Query features + model-pool accuracy/energy history & No & No & Yes (online bandit) & Still no cloud queue, RTT, SLA or edge-health channel \\
NSGA-II router (Yu et al., 2025) & Offline Pareto front over quality/latency/cost & Yes, offline & No & No & Fixed once optimised, so cannot react to a live queue or RTT spike \\
PROTEUS (Bhatti et al., 2026) & Accuracy target $\tau$ via Lagrangian RL & No & No (accuracy floor, not a deadline) & No & State space excludes queue, latency and energy entirely \\
Splitwise (Patel et al., 2024) & Phase-level cluster telemetry & Yes, cluster-level & No & No & Scheduling is cluster/phase granularity, not a per-request deadline \\
LiteLLM (LiteLLM, 2026) & Weighted/\-latency/\-usage/\-least-busy rule, operator-selected & Reactive only (latency cache, cooldown) & No & No & Fixed rules over a rolling window, not a learned or self-reported signal \\
Portkey (Portkey, 2026) & Static operator-set traffic weights & No & No & No & Weights are configured up front, not adapted from live telemetry or outcomes \\
\textbf{This work} & 13-slot masked state: confidence, cloud queue, RTT, SLA remaining, edge resource stress, calibration delta, error rate, operational state & Yes, 5 live signals + edge self-report & Partial --- per-request signal, learned soft preference via reward, not a hard gate & Yes, bidirectional loop & Still over-escalates in scenarios where RTT alone already guarantees an SLA violation (Section 7.6) \\
\bottomrule
\end{longtable}
}

Three gaps recur across every system in Table~\ref{tab:comparison},
and each maps directly onto a specific component of the design in
Section 5, rather than being a claim added after the literature was
read. First, none of the confidence- or budget-only cascades
(BranchyNet, MSDNet, EdgeBERT, SPINN) nor any learned or production
router (RouteLLM, GreenServ, LiteLLM, Portkey) reads live infrastructure
state at all: the Edge Context
Protocol and the Bronze/Silver/Gold Medallion pipeline exist
specifically to give the routing decision cloud queue depth, RTT, SLA
remaining and edge resource stress alongside confidence, rather than
confidence alone. Second, the offline-optimised NSGA-II router and the
accuracy-floor PROTEUS both freeze a policy that, once trained, cannot
see the request-time queue or RTT. The PPO agent in this project is also
trained offline, but its state and reward both carry the current
request's own SLA remaining and predicted RTT, at request-time
granularity. None of the eleven systems above matches this. Splitwise
leaves it at cluster granularity instead. This is a learned
preference the policy weighs against everything else, not a
deterministic per-request gate --- Section 5 discloses this distinction
directly, including a Silver-layer signal computed for exactly a harder
guarantee but not yet wired into any decision path, so the honest claim
here is request-time granularity, not a guarantee no surveyed system
offers. Third,
none of the eleven systems closes the loop back to the requester: EdgeBERT
adapts within a request and GreenServ adapts its bandit online, LiteLLM
adapts its latency cache and cooldown window, but none of them feeds the
outcome of one routing decision back into the state
the next decision is made from at the level of an individual device.
The Bidirectional Loop (Section 5) is designed against exactly that
gap, and its removal is measured directly, not only asserted, in
Ablation Run 4 (Section 7.3).

